\scanpage{225}
\mspage{141}
% archive sheet number: 224
% element: text | top, continuation of transfinite construction | le premier ordinal
le premier ordinal tel que $\operatorname{card}\alpha_0>\pi_0$, \inserted{donc l'ens $\bar I_{\alpha_0}$ des ord. $\alpha<\alpha_0$ est grd devant $\pi_0$.} Je dis \deleted{que} qu'on a $C_{\alpha_0}=C$. Pour ceci, il suffit de voir que $C_{\alpha_0}$ est stable par $\unsure{\underline{\lim}|\varinjlim}$ de cardinal $\leq\pi$. Mais si $I$ est une catégorie telle que $\operatorname{card}\mathrm{Fl}\,I\leq\pi_0$,
% element: diagram | upper third, first functor and inserted union | phi I to C alpha zero
\diagram{\begin{tikzcd}[ampersand replacement=\&]
\unsure{\varphi|\phi}:\ I\arrow[r]\&C_{\alpha_0}\inserted{=\bigcup_{\unsure{\alpha'|\alpha^\ast|\alpha^\sharp}<\alpha_0}C_{\unsure{\alpha'|\beta'}}}
\end{tikzcd}}
% element: text | upper middle, factorization | foncteur alors
foncteur, alors la définition de $\alpha_0$ implique que $\varphi$ factorise en
% element: diagram | middle, factorization plain stroke | phi alpha I line C alpha prime
\diagram{\begin{tikzcd}[ampersand replacement=\&]
\varphi_\alpha:\ I\arrow[r,dash]\&C_{\alpha'},
\end{tikzcd}}
% element: margin note | left middle, slanted question and correction | Un peu court
% slanted
% The last two lines are crossed; the large crossing also intersects earlier strokes. The short word after est and the cancelled opening remain uncertain.
\margin{? Un peu court. Il faut invoquer le fait que \deleted{\unsure{card}} $\bar I_{\alpha_0}$ est \unsure{le|b.} grd devant $\pi_0$\unsure{;|,} \deleted{(\unsure{Il|$\bar I_{\alpha_0}$} \unsure{est} \unsure{$=$|$\bar m$} grd devant \unsure{$\forall$|tt} $\unsure{\pi'_0|\pi_0}$-\unsure{quelconque})}}
% element: text | middle, closure conclusion | pour alpha prime less than alpha zero
pour $\alpha'<\alpha_0$. Donc $\underline{\lim}\varphi_{\deleted{\unsure{\alpha_0|\alpha'}}}\in\mathrm{Ob}\,C_{\alpha+1}\subset\mathrm{Ob}\,C_{\alpha_0}$, ce qui montre la stabilité requise.

% element: text | middle lower, essential smallness | Pour voir que C est equiv
Pour voir que $C$ est équiv. à une petite catégorie, il suffit de le voir pour les $C_\alpha$, \unsure{par|en} procédant par induction transfinie. Cela résulte du lemme suivant, immédiat:

% element: text | lower third, lemma 4.6.2 | Lemme La limite inductive
% The initial article is heavily retraced; La, an overwritten Toute, and underline/cancellation states are explicitly alternative readings.
\uline{Lemme} 4.6.\unsure{1|\deleted{1}}2. \unsure{1°)|a)} \unsure{\uline{La}|\deleted{Toute}\uline{La}|\uline{Toute}|\deleted{La}} limite inductive \inserted{$C$} d'un syst. inductif filtrant $(C_i)_{i\in I}$ de catégories $C_i$ équivalentes à des petites catégories, est équivalente à une petite catégorie. Si $\unsure{\operatorname{card\,rél}|\operatorname{card\,réd}}(C_i)\leq\pi$ \inserted{pour $\forall i$,} et $\operatorname{card}I\leq\pi$, alors $\operatorname{card\,red}C\leq\pi$ ($\pi$ étant un
